\documentclass{WileyMSP-template}


%++++++++++++++++++++++++++ TEMPORARY SECTION, DELETE BEFORE SUBMISSION +++++++++++++++++++++++++++++++++++
\usepackage[dvipsnames]{xcolor}
\usepackage{amssymb}
%\usepackage{natbib}
\newcommand{\jb}[1]{\textcolor{blue}{#1}}
\newcommand{\np}[1]{\textcolor{ForestGreen}{#1}}
\newcommand{\ds}[1]{\textcolor{purple}{#1}}
\newcommand{\ik}[1]{\textcolor{RawSienna}{#1}}
\usepackage{lineno}
\linenumbers
\usepackage{fancyhdr}
\usepackage{lastpage}

%added by np
\usepackage[normalem]{ulem}
%++++++++++++++++++++++++++++++++++++++++++++++++++++++++++++++++++++++++++++++++++++++++++++++++++++++++++

\begin{document}


\pagestyle{fancy}
\rhead{\includegraphics[width=2.5cm]{vch-logo}}


\title{Degradation at Work: Probing Irreversible Processes in Organic Radical Batteries by Electron Paramagnetic Resonance Spectroscopy}

\maketitle



% Author: Please give full first and last names for authors and include * after the name of all corresponding authors

\author{Ilia R.\ Kulikov}%\thanks{These authors contributed equally} #np this needs to be in acknowledgements not here

% Dedication

\dedication{}



% Affiliations: Please provide adacemic titles (Prof. or Dr.) for all authors where applicable, and include an institutional email address for all corresponding authors

%\affiliation{$^1$Berlin Joint EPR Lab, Freie Universit\"{a}t Berlin, Fachbereich Physik, Arnimallee 14, D-14195 Berlin, Germany}
%\affiliation{$^2$Electrochemistry Department, St.\ Petersburg State University, 7/9 Universitetskaya nab., St.\ Petersburg, 199034, Russian Federation}


\begin{affiliations}
I. R. Kulikov\\
Berlin Joint EPR Lab, Freie Universit\"{a}t Berlin, Fachbereich Physik, Arnimallee 14, D-14195 Berlin, Germany\\
Email Address: j.behrends@fu-berlin.de\\

\end{affiliations}


% Keywords: Please provide a minimum of three and a maximum of seven keywords, separated by commas

\keywords{EPR; Degradation; Organic Radical Batteries; In-operando; Organic Cathode Materials}


% Abstract should be written in the present tense and impersonal style (i.e., avoid we), and be at most 200 words long
\begin{abstract}

\ik{
Using the methods of electron paramagnetic resonance (EPR) spectroscopy and electro-chemistry, we observed a change in the number of charge-bearing groups in a nitroxide-containing polymer cathode of an organic radical battery. We present a design of an electrochemical cell suitable for room temperature, operando EPR characterization during repeated charge-discharge cycling. The two-component deconvolution of the EPR spectra allowed us to identify the paramagnetic species that leave the polymer film during the cycling. We performed a quantitative analysis of the EPR spectra to show that the  electrochemical cell based on the redox-conductive polymer loses its capacity in its oxidized, or charged state.\\
We further studied the polymer cathode films of the cells at cryogenic temperatures and distinguished between the subsystems of weakly and strongly interacting charge-bearing nitroxide fragments that allowed us to propose two possible mechanisms of the degradation.
}

\end{abstract}

% Text: Please use section headings and subheadings as specified below. For communications, all section headings apart from Experimental Section should be removed
% Please make the first reference to a display item bold: \textbf{Figure 1}
% Do not abbreviate Figure, Equation, etc.; display items are always singular, i.e., Figure 1 and 2.
% Equations are always singular, i.e., Equation 1 and 2, and should be inserted using the {equation} environment, not as graphics
% Please do not use footnotes in the text, additional information can be added to the Reference list.


\section{Introduction}
Batteries based on conjugated polymers containing stable radical substituents as high-capacitance groups represent a promising class of future electrochemical power sources.\cite{nakahara2002_cpl, nishide2004_electact,xie2021_mathoriz} They combine the advantages of high-power supercapacitors, namely high discharge rates, and the high energy density of conventional lithium-ion technology. A further beneficial property of organic materials over traditional inorganic materials is their availability and the low cost of the starting materials for the synthesis of the target polymers in conjunction with good mechanical properties.\cite{janoschka2012_advmater, muench2016_chemrev, friebe2017_topcurrchem}
\bigbreak
While active electrode materials with nitroxide radicals as redox-active groups are ideally suited for organic radical batteries (ORBs) that exhibit high power densities, the broad application of most nitroxide-based materials is limited by their moderate electrical properties. A promising route towards overcoming the conductivity problem is the use of polymers that combine radical-containing moieties and a conductive backbone. This strategy was successfully followed in a number of studies focusing on different polymers.\cite{oyaizu2015_polymerjournal, bahaceci2013_jpowersources, katsumata2006_mrc, xu2014_electact, aydin2015_jsoistatelect, schwartz2018_synthmet} 
\bigbreak
Common to all these polymers is the fact that they contain 2,2,6,6-tetra-methylpiperidine-1-oxyl (TEMPO) as nitroxide pendant groups. Upon changing the redox state of the TEMPO group, the neutral nitroxide radial is converted into the aminoxyl anion or the oxoammonium cation. Only the neutral radical has an unpaired electron spin ($S = 1/2$) and can thus be detected using Electron Paramagnetic Resonance (EPR) spectroscopy. For this reason EPR is commonly used to detect the presence of radicals in active-electrode materials for ORBs.\cite{nakahara2002_cpl, nishide2004_electact, bahaceci2013_jpowersources, aydin2015_jsoistatelect, khodeir2019_softmatter} The vast majority of EPR studies use continuous wave EPR (cwEPR) measurements merely to confirm the radical nature of the active materials ex-situ. However, also in-situ and in-operando experiments have recently been demonstrated, providing direct access to the change of the redox state upon charging and discharging ORBs.\cite{huang2016_jpowersources, kanzaki2018_acsappmat} This opens the intriguing perspective of determining yields for charging and discharging processes by means of in-operando quantitative cwEPR spectroscopy. We note that quantitative cwEPR can accurately assess the absolute number of unpaired electron spins in a given volume with high sensitivity. This has, for instance, been recently used to quantitatively investigate the doping of conjugated polymers.\cite{arvind2020_jpcb, pingel2016_aem}
\bigbreak
Conventional in-operando cwEPR experiments were previously employed to monitor structural changes at the anode in lithium-ion batteries,\cite{wandt2015_ees} highlighting the usefulness of in-operando EPR spectroscopy for battery research in general. The cwEPR spectrum itself contains valuable information on the molecular environment and the local geometry of the spin-bearing groups. For instance, in the case of nitroxide-based materials, the cwEPR spectra of immobilised dilute nitroxides in a solid film, densely packed nitroxides in a thin film, nitroxides in solution and two nitroxides in close vicinity forming a dimer are all well distinguishable.\cite{santangelo2008_jpcb,clark2016_concmagnres}
\bigbreak
However, in-operando experiments are by no means restricted to cwEPR measurents. The information content accessible by EPR can substantially be increased further by using modern pulse EPR (pEPR) techniques.\cite{schweiger2001_book} This is particularly true for ORBs with TEMPO-based active electrode materials, as nitroxide radicals give strong and well-studied EPR signals. In fact, nitroxide radicals are probably the most studied paramagnetic system due to their routine use in EPR-based distance measurements\\ in proteins.\cite{jeschke2012_annrevphyschem} Apart from accurately determining distances between neighbouring radicals, pEPR techniques can also provide new insights into the wave function localisation of paramagnetic states~\cite{steyrleuthner2017_pccp} as well as molecular motion of the nitroxide groups.\cite{clark2016_concmagnres}
\bigbreak
Here we describe the development of an electrochemical cell that allows us to perform in-operando cwEPR measurements as well as pEPR experiments using conventional EPR resonators. We employ quantitative cwEPR spectroscopy to study charging and discharging processes of a novel, recently synthesized electrode material for ORBs based on a di-TEMPO Ni-Salen polymer (p-DiTS).\cite{vereshchagin2020_battandsupercaps} Our results demonstrate that almost all TEMPO groups can successfully be oxidised and reduced upon changing the voltage applied to the cell. The distinct differences between the room-temperature EPR spectra of nitroxide radicals immobilised on the cathode and those dissolved in the electrolyte enable us to explore the influence of various deposition parameters on the stability of the p-DiTS film prepared by electropolymerisation. Further, we adjust the redox state of p-DiTS in a cwEPR experiment by choosing the applied electrical potential accordingly. The desired state of the cell is then preserved by freezing using cryogenic helium. The sample is subsequently transferred to a pEPR setup. We successfully conduct electron spin echo measurements that provide the basis for electron-electron double resonance (ELDOR) experiments for accurately determining distances between adjacent paramagnetic centres.

\bigbreak
\subsection*{Versatile Electrode Setup for ex-situ and in-situ Spectroelectrochemical EPR}

\np{
There is increasing interest in both ex-situ and in-situ spectroelectrochemical (SEC) EPR measurements in research fields ranging from biology to materials science, for redox active systems such as proteins and hydrogenase to metal catalysis, redox-flow batteries and of course ORBs. One significant drawback with SEC EPR is the fact that microwave (mw) resonators suffer from substantial mw damping as a result of introducing the metal electrodes, polar solvent and ionic salts (i.e., the electrolyte) needed for successful electrochemistry. The damping of the mw can be significantly reduced by our novel on-substrate electrode design (see Figure \ref{fig:Figure_1}) as this limits the amount of metal that is introduced into the resonator, while still allowing us to have electrode surface areas large enough to deposit sufficient material of interest to observe EPR signals.
}


\begin{figure}
	\includegraphics[width=0.96\textwidth]{Figure_1}
	\caption{Schematic diagram of the electrode design and assembly of the electrochemical cell based on a modified X-band EPR tube. \textbf{a)} Molecular structure of p-DiTS, with the redox reaction of the charge-bearing TEMPO groups; Highlighted is the backbone, p-NiMeOSalen. \textbf{b)} Setup for deposition and cleaning of polymer films. \textbf{c)} Tube-based electrochemical setup. d) Photograph of the modified quartz EPR tube with a flattened bottom section, and the process of forming the modified tube. \textbf{e)} Cyclic voltammograms of a growing p-DiTS film during 100 electrochemical polymerisation cycles (solid), during the cleaning process (crossed) and in the tube-based setup (dotted). Peaks A/A', B/B' correspond to oxidation and reduction of the TEMPO fragments in p-DiTS and the NiSalen backbone, respectively. All CV were recorded at 50 mV/s.}
	\label{fig:Figure_1}
\end{figure}

\bigbreak

\np{	
The on-substrate electrodes were produced as follows. Firstly, we use polished quartz substrates, which do not give EPR signals, as our base. We clean these quartz substrates by sonicating in Acetonitrile (ACN) followed by Isopropyl alcohol and finally Ethanol, before placing in a UV Ozone cleaner (Ossila) for $\sim$10~minutes. Clean substrates are then placed inside a holder with shadow masks ready for thermal evaporation and positioned inside a thermal evaporation chamber which is integrated into a Nitrogen glovebox (MBraun ProVap 5G PVD System equipped with a thermal evaporator). The evaporation chamber is pumped to $10^{-6}$~mbar after which we evaporate 10~nm of Chromium (Cr) as an adhesion layer followed by 180~nm of Gold (Au) as the working and counter electrodes (WE and CE respectively). The reference electrode (RE) we do not evaporate on-substrate as we find stability of the cyclic voltammograms (CVs) is better with a $\varnothing$250~$\mu$m Ag wire which is either used as is or for additional stability coated galvanically with a AgCl layer (see SI \ref{Ref_electrode} for details of coating Ag with a AgCl layer). 
}

\bigbreak

\np{
The above procedure allows a WE active area of 12.0~mm$^2$ which allows us to deposit enough material of interest for the current study. Our versatile on-substrates electrodes can easily and readily be modified for different geometries and of course electrode active areas, by simple modification of the shadow mask used. Two examples that show the merit of our flexible design are as follows: 1) we can increase the WE area for samples with particularly small EPR signals while still maintaining particular film thicknesses of interest (allows study of EPR properties as a function of film thickness). 2) In cases where the electrochemical process is thought to deposit material on both the WE and CE we can control the distance between the two electrodes and their relative position on the substrate, where for example moving the CE further away from the most sensitive area of the microwave cavity would allow selective probing of the WE and vice versa for the CE.
}

\bigbreak

\np{
While the on-substrate electrodes reduced microwave damping, this is not the only cause of the microwave damping at room temperature. A substantial amount of microwave damping occurs due to the electrolyte, made of polar solvent and ionic salts. Therefore, we further optimize our setup by modifying existing 5~mm outer diameter (OD) Quartz EPR tubes by flattening the bottom 1.3$-$1.6~cm (cf. Figure~\ref{fig:Figure_1}b), thereby reducing the active electrolyte needed to submerge both the WE and CE on the substrate and of course the RE wire, from $\sim$120~$\mu$L to $\sim$60$-$70~$\mu$L. 
}

\bigbreak

\np{
The flattened tube and flat electrode setup also have an additional benefit in that it moves the sample further away from the electric field distribution in the Bruker ER~4122-SHQE resonator (TE\textsubscript{011} mode cavity for cwEPR). Thereby further increasing the resonator quality factor (Q-factor, given by $Q = \frac{f_{res}}{\Delta f}$, $f_{res} =$ Resonance frequency, $\Delta f =$ FWHM of resonance dip) and hence sensitivity. Further improvements can be made using our setup in a cylindrical cavity such as a Bruker ER~4103-TM (TM\textsubscript{110} mode cavity for cwEPR, developed for studying samples exhibiting high dielectric constants), where the flattened cell can be aligned to further reduce coupling to the microwave electric field. Our modified tube is essentially a home-made ``flat cell'', with the additional important benefit that our modified tubes are compatible with commercial Bruker ER~4118X-MD5 resonators, most commonly used for advanced pulse EPR measurements at X-band frequencies ($f = 9-10$~GHz).
}

\bigbreak

\subsection*{Materials}

\ik{The materials investigated in this study are 1) p-DiTS - a redox-conductive polymer with a NiSalen conductive, conjugated backbone and two charge \np{bearing} \sout{baring} redox-active nitroxyl radicals (TEMPO) per monomer unit, 2) polymeric NiMeOSalen (p-NiMeOSalen) which is chemically similar to p-DiTS except it does not contain the TEMPO-succinyl moieties and instead is terminated with a methyl group after the remaining oxygen, see Figure \ref{fig:Figure_1}a for both molecular structures.}

\bigbreak
%\ik{P-DiTS can be anodically electropolymerized from the monomer in an electrolyte solution into thin films, representing the active electrode materials in ORBs. The TEMPO groups in p-DiTS can undergo a reversible and very rapid redox reaction between two stable states, a neutral nitroxyl radical (TEMPO$^{\bullet}$) and an oxoammonium cation (TEMPO$^{+}$), that are shown in Figure \ref{fig:Figure_1}~a). The conductive backbone together with the high concentration of the unpaired electrons endow p-DiTS films with high specific capacity and fast charge and discharge rates (83$\,$mAh$\,$g$^{-1}$ at 60$\,$C and 31$\,$mAh$\,$g$^{-1}$ at 3000$\,$C, with 66\% capacity retention after 2000 charge-discharge cycles)\cite{vereshchagin2020_battandsupercaps}. That makes p-DiTS a promising cathode-active charge storage material.
%}

\np{P-DiTS can be anodically electropolymerized from the monomer in electrolyte solution to an on electrode thin film -- representing the active cathode material in a ORB. The NiMeOSalen conductive backbone together with the high concentration of unpaired electrons (TEMPO radicals) endow p-DiTS films with high specific capacity and fast charge and discharge rates (83$\,$mAh$\,$g$^{-1}$ at 60$\,$C and 31$\,$mAh$\,$g$^{-1}$ at 3000$\,$C, with 66\% capacity retention after 2000 charge-discharge cycles)\cite{vereshchagin2020_battandsupercaps}. This makes p-DiTS a promising cathode-active charge storage material.}

\bigbreak

\np{The TEMPO groups in p-DiTS can undergo a reversible oxidation and rapid redox reaction between two stable states~\cite{vereshchagin2020_battandsupercaps}, a neutral nitroxyl radical (TEMPO$^{\bullet}$) and an oxoammonium cation (TEMPO$^{+}$), see Figure \ref{fig:Figure_1}a. In the reduced (discharged) state, the TEMPO$^{\bullet}$ moieties of p-DiTS possess unpaired electrons that makes them paramagnetic and therefore EPR active. In the oxidized (charged) state, the loss of an electron results in TEMPO$^{+}$ which is diamagnetic and EPR silent. Polymeric NiSalen (e.g., p-NiMeOSalen) is also known to be redox active and to show potential-dependent EPR signals, which were attributed to biphenoxyl radicals by Dmitrieva \textit{et al}~\cite{Dmitrieva2018}. Therefore we study both p-DiTS and p-NiMeOSalen using SEC cwEPR.
}

%\bigbreak

%Threfore, a cwEPR spectrum for p-DiTS may contain not only the signal from the TEMPO radicals, but also a contribution from the positive polarons on the oxidized NiSalen backbone. In order to distinguish a possible spectral contribution of the oxidized backbone in a potential-dependent p-DiTS spectrum, films of poly-Ni-Methoxy-Salen (p-NiMeOS, Figure \ref{fig:Figure_1}, inset) were prepared and their potential-dependent in-situ and ex-situ cwEPR spectra were recorded and compared to the spectra of p-DiTS.\\

\bigbreak

\subsection*{Electrochemical Polymerization and Cyclic Voltammograms in a Modified EPR Tube}\label{CVTube}

\np{
The first step of the sample preparation for SEC EPR is the electropolymerization of monomeric DiTS to p-DiTS and simultaneous deposition onto the on-substrate Au WE. To do this we use the setup shown in Figure~\ref{fig:Figure_1}c, where we use the on-substrate 12~mm$^2$ Au electrode (shown in Figure~\ref{fig:Figure_1}) as the WE and a coiled Pt wire for the CE which are submerged in the electrolyte (0.1~M Et\textsubscript{4}NBF\textsubscript{4} in ACN) containing 1~mM monomeric DiTS. The chemically inert RE used was an Ag wire in AgNO\textsubscript{3} solution which was placed inside a glass tube with a porous CoralPor\textsuperscript{TM} tip; placed inside the electrolyte solution. To carry out the electropolymerization and deposition onto the on-substrate Au WE, we apply between 8 and 200 CV cycles from $-$200/$-$50~mV to $+$900~mV at a rate of 25 or 50~mV/s. In studies involving NiMeOSalen we apply between 10$-$20 CV cycles for electropolymerization/deposition. The same number of CV cycles produces thicker p-NiMeOSalen films compared to p-DiTS films as in DiTS the coulombic repulsion provided by the oxidized TEMPO groups hinder film growth after $\approx$200 CV cycles~\cite{vereshchagin2020_battandsupercaps}. We estimate the thickness ($t$) of the p-DiTS films used in this study to be the following, $t\approx$40~nm for 8~CV cycles, $t\approx$500~nm for 200~CV cycles; and for p-NiMeOSalen $t\approx$900~nm for 15~CV cycles, see SI \ref{fig:S_SAMPE_INFO} for details regarding thickness estimation.
}

\bigbreak

The most electrochemically stable films came from polymerized systems which were cleaned by rinsing in clean ACN followed by Propylene Carbonate (PC) and lastly CV cycling in clean electrolyte. The effects of rinsing and cycling of the polymer films in clean solvents/electrolyte are described later on.

\bigbreak

%\ik{The flattened region of the modified EPR tube has a greatly decreased volume of the electrolyte as compared to the deposition and the cleaning beakers and the standard EPR tube. There is only 0.35~mm space between the substrate and the flattened wall of the tube that is shared between the front and the back side of the flat, modified region. That leaves only an $\approx200$~um layer of the electrolyte between the WE and the CE. For the cells in the modified tubes, the CV had to be taken at lower cycling speeds, otherwise the peak widths and positions were distorted and/or the voltammograms had an arbitrary linear shift in the direction of higher voltages in the range of 100--500~mV.\\
%An example of a CV with slightly distorted peaks for a modified tube is shown in Figure \ref{fig:Figure_1}~e) in a dotted blue line.
%}

\np{While we see a significant improvement in the cwEPR by using our modified ``flat cell'' EPR tube, there are some difference in the CV curves measured in a beaker (setup schematic Figure \ref{fig:Figure_1}c) vs. in a modified tube (setup schematic Figure \ref{fig:Figure_1}d). As seen in Figure \ref{fig:Figure_1}e the CV measured in our modified tube (blue curve) shows slightly shifted and broader oxidation and reduction peaks as compared to the in beaker measurements (gold and black curves) for the same p-DiTS film. We attribute this effect to the smaller electrolyte volume used in the modified tube. The electrolyte volume seen by the on-substrate films in the modified tube is even less then the total electrolyte volume of 60 -- 70~$\mu$L, as only 0.35~mm of space remains between the inner width of the flattened modified tube and the substrate. Hence, we estimate that no more then a 0.2~mm layer of electrolyte is seen by the films. The CV in a modified tube show less distortions with slower CV scan rates (e.g. 5~mV/s used for studies in section \ref{Degradation_study}) as this gives more time for circulation of electrolyte in the tube. We note that while there is some distortion, this is relatively minimal and importantly the oxidation and reduction peaks remain rather symmetric suggesting that the reversibility of the redox process is not hindered.} 
\bigbreak

\subsection{EPR-Active States in pDiTS-based Electrochemical Cells}\label{Classification}
\label{subs:classification}

\bigbreak

EPR spectra of electrochemical cells with pDiTS as active-electrode material can exhibit a variety of different signals. Before presenting and discussing the experimental results, we will provide an overview about the various characteristic cwEPR lineshapes associated with nitroxides in different environments.

\bigbreak

\textbf{Figure~\ref{fig:Figure_2}}a shows the spectrum expected for a \emph{\textbf{densely packed pDiTS film}} with all nitroxides in the EPR-active reduced state. The broad unstructured line results from the strong dipolar and exchange interaction between both nitroxides of each DiTS monomer as well as between the nitroxides of different monomers. \np{The exchange interaction can cause both line narrowing and broadening depending on the strength of the interaction, while the dipole interaction causes broadening, this means the overall cwEPR line width depends on the spin concentration in the film.}% Dipolar and exchange interactions can cause both line narrowing and line broadening, resulting in the fact that the overall line width depends on the spin concentration in the film.

\bigbreak
\begin{figure}
	\includegraphics[width=1\textwidth]{Figure_2}
	\caption{Classification of the nitroxide cwEPR spectra observed for p-DiTS under different conditions and environments. a) Densely packed nitroxide film with exchange and dipolar broadening or exchange narrowing giving a single derivative, b) typical dilute mono-nitroxide in solution spectrum, three-lines due to $S=1/2$ $I=1$ $^{14}N$ hyperfine interaction, c) typical di-nitroxide in solution where dynamic changes in nitroxide distance modulated the exchange interaction giving five lines (actual cwEPR of DiTS monomers in solution at room temperature), d) Immobilized nitroxide spectrum with both $g$ and $A$ (Hyperfine) anisotropy.}
	\label{fig:Figure_2}
\end{figure}

\bigbreak

\emph{\textbf{Dilute mono-TEMPO fragments}} dissolved in the electrolyte, each containing only one nitroxide, give rise to a three-line spectrum (cf.\ Figure~\ref{fig:Figure_2}b) which is well known from nitroxide-based spin labels in solution. Tumbling of the molecules in the electrolyte results in an averaging of the anisotropic interaction between the electron spin and the nitrogen nuclear spin ($I = 1$) as well as of the $g$ anisotropy, leaving only the isotropic part of the hyperfine coupling and the isotropic $g$ value. 

\bigbreak

\emph{\textbf{Dilute DiTS monomers}} in the electrolyte generally give a spectrum consisting of five narrow lines (cf.\ Figure~\ref{fig:Figure_2}c), if each monomer contains two interacting EPR-active nitroxides. The spectrum depends on dynamic effects, specifically the modulation of the exchange interaction between both radicals of each monomer.  Thus, depending on the solvent and temperature, the spectrum can also rather look like the three-line spectrum of mono-TEMPO fragments.

\bigbreak

The fully charged (oxidized) pDiTS film can contain domains in which the nitroxides are not connected to the electrode and cannot be oxidized or reduced. These electrically inactive \emph{\textbf{immobilized DiTS \\monomers}} in the active-electrode film display a broad spectrum as shown in Figure~\ref{fig:Figure_2}d. In contrast to the densely packed pDiTS film, there is no substantial exchange interaction between neighboring radicals. DiTS monomers with either one or two paramagnetic nitroxides result in similar EPR spectra.

%The oxidized TEMPO fragments in the electrolyte, the oxidized pDiTS fragments inside the film and other diamagnetic species do not exhibit spin resonance and are EPR silent. \np{do we need this here - the section is on EPR active signatures.}

\bigbreak

We note that the oxidized NiSalen backbone can contribute to the spectrum as well. Its EPR signature, which is not shown in Figure~\ref{fig:Figure_2}, is clearly different from the \np{nitroxide/TEMPO}-related spectra discussed above.


\section{Results and Discussion}

\np{
To establish the extent to which our modified tube makes a difference to both the resonator Q-factor and the signal to noise ratio (S/N), we measure the cwEPR spectrum of p-DiTS deposited on our electrode setup with the tube filled with electrolyte to a height which submerged both the WE and CE, this requires 100~$\mu$L in the standard 5~mm OD quartz EPR tube and 70$\mu$L in our modified tube, and also record the Q-factor associated with each spectrum as a function of height away from the resonator centre $H\textsubscript{res. center}$.}

\begin{figure}[h!]
	\includegraphics[width=0.5\textwidth]{Figure_3}
	\caption{Changes in the Q factor of a microwave resonator loaded with a p-DiTS electrochemical cell with respect to the distance between the center of the red-ox active cathode and the center of the resonator. A standard 5$\,$mm OD tube and the modified tube with the flat tip. Inset: X-band cwEPR spectra for the p-DiTS cell for the standard tube and for the modified tube inserted at $H\textsubscript{res. center}$= 19~mm, recorded with the 0.5~mT modulation amplitude at a microwave power of 1.26~mW.}
	\label{fig:Figure_3}
\end{figure}
\bigbreak

\np{We determine that an empty resonator has Q$\approx$7800, while  inserting the tubes at $H\textsubscript{res. center} =$ 25~mm (with on-substrate electrodes and electrolyte) reduces the Q for both cases, but more significantly for the standard tube $Q\textsubscript{standard tube}\approx$6000, $Q\textsubscript{modified tube}\approx$7600, see Figure \ref{fig:Figure_3}. At $H\textsubscript{res. center}=$19~mm the Q-factor drops for both, again more significantly for the standard tube to $Q\textsubscript{standard tube}$ $\approx$630 while for the modified tube it is still $Q\textsubscript{modified tube}$$\approx$3000, a factor of $\approx$5 difference. This a significant difference between the two tubes and shows the strength of using our modified tube and flat electrode geometries for SEC EPR. Furthermore, inserting the modified tube to only $H\textsubscript{res. center}=$12~mm, we still have a Q-factor of $Q\textsubscript{modified tube}$$\approx$850 which is still higher then in the standard tube at a $H\textsubscript{res. center}=$19~mm. Coupling the resonator with the standard tube setup was not possible below $H\textsubscript{res. center}=$19~mm.  
}

\bigbreak

\np{
The decreasing Q-factor is not the only affect of a high-dielectric sample in a EPR resonator. The interaction of the high-dielectric sample with the electric field distribution of the resonator causes observation of a mixture of the absorption and dispersion signals, leading to asymmetric lineshapes in the cwEPR spectra. This makes double integration of the derivative cwEPR signal more complex, thereby making quantitative analysis more difficult. This affect can be seen in the cwEPR spectra (Figure \ref{fig:Figure_3} inset) of p-DiTS measured in the standard 5~mm tube where the cwEPR signal is asymmetric, but this is not the case when the modified tube is used where we see symmetric lineshape at the same sample height $H\textsubscript{res. center} = 19$mm, further details are presented in SI~\ref{Q-factor_SNR}. 
}

\bigbreak

\np{Furthermore, the significance of having the cell inserted as close to the resonator center as possible is shown by the factor of $\approx$2 increase in the S/N, when comparing the EPR spectra measured in the modified tube and standard tube at maximum sample insertion. This is most likely a lower estimate of the increased S/N ratio as the signal intensity was measured at the peak intensity, which is distorted for the standard tube due to a mixture of absorption and dispersion signals. }

\bigbreak

\np{
Using the modified tube we can achieve higher Q-factors and therefore higher sensitivity; increased the S/N as compared to measuring in a standard EPR tube and also symmetric cwEPR lineshapes which allow for more straightforward quantitative analysis of the signals, especially useful for in-situ cwEPR with varying redox potentials. 
}
\bigbreak

\subsection*{p-DiTS \textit{In-Operando} Redox Series}

\np{Potential dependent in-operando SEC cwEPR was measured on a 40~nm thick (8~CV cycle) p-DiTS film, inside our modified EPR tube at room temperature (see Figure \ref{fig:Figure_4}). At 0.3~V (vs. Ag) the p-DiTS film is in the TEMPO$^{\bullet}$ reduced/discharged state. Here we see a single broad derivative lineshape, with a peak to peak line width of ....mT, which is due to the high spin concentration in the film where the line width is due to a combination of dipolar and exchange coupling. We fined that in this reduced neutral radical state the number unpaired electrons/ spins in the film is $6.4\times10^{14}$ (up to 20\% error). This is comparable to the number of electrochemically active TEMPO charges at $7.6\times10^{14}$, determined by integration of the cathodic branch of the CV of the cell vs. time.}

\bigbreak

\np{Upon increasing the oxidation potentiality, we observe no significant change in the spectral shape of the acquired cwEPR spectra in the oxidation potential range of 0.3~V to $\sim$0.75~V (region I). A representative cwEPR spectrum from region I, taken at oxidation potential 0.6~V is shown in Figure \ref{fig:Figure_4}c. We also see no significant change in the number of spins (spin count) contributing to the cwEPR spectra at all potentials in region I, as shown in Figure \ref{fig:Figure_4}b. Increasing the oxidation potential further to 0.8~V starts to lead to a decreasing spin count, and from 0.85~V to 0.94~V (the maximum applied for this cell in-operando) there is a sharpe and significant drop in the spin count. We observe the smallest spin count at an oxidation potential of 0.94~V giving $5.4\times10^{13}$, an order of magnitude decrease in the number of unpaired electrons in the film compared to the initial state.}

\bigbreak 

\begin{figure}[t!]
	\includegraphics[width=1\textwidth]{Figure_4}
	\caption{Room temperature X-Band (9.40$\,$GHz) cwEPR spectra for a p-DiTS electrochemical cell under operating conditions. a): cwEPR spectra during charging (oxidation) and discharging (reduction) of the cell. b): number of unpaired electron spins determined from the quantitative analysis of the spectra (EPR spin count) and number of electrochemically active charges in the cell determined from its cyclic voltammogram (CV). c): slices for the initially discharged state ('I', 0.6$\,$V), charged state ('II', 0.90$-$0.94$-$0.82$\,$V, averaged within the region denoted by $\updownarrow$~, scaled up $\times$10), and the final discharged state ('III', 0.6$\,$V).}
	\label{fig:Figure_4}
\end{figure}

\bigbreak

\np{Upon closer inspection, by averaging the cwEPR spectra in region II (0.90-0.94-0.82~V) Figure \ref{fig:Figure_4}c, we clearly see the spectral shape has substantially changed from the initial region I spectrum. The region II averaged spectrum can be decomposed using simulations via the MATLAB toolbox easyspin and the pepper function, into two spectral components. The first component is a broad derivative with a linewidth of ....mT (Figure \ref{fig:Figure_4}c orange trace), similar to that of region I and a second component consisting of a 3-line spectrum.}

\bigbreak

\np{The presence of component 1 even at the oxidation potential of 0.94~V is at first unexpected, as at this potential the CV curve shows full TEMPO oxidation. We therefore suggest that perhaps  component 1 corresponds to local islands of still densely "packed" spins, which are electrically disconnected from the electrodes and hence not redox active. This idea will be further explored later in the text.}

\bigbreak

\np{The presence of component 2 (3-line feature) is as discussed in section \ref{Classification} due to mobile nitroxide spins. The presence of a 3-line component but no detectable 5-line component suggests the nitroxide spins in electrolyte are not due to release of monomers (as expected for this film as it under went the cleaning procedure described in section \ref{CVTube} to remove any non-polymerized DiTS. The 3-line feature must therefore be due to molecular breakdown of DiTS to give mono-nitroxide moieties in solution.}

\np{The spin count for component 1 ``packed'' immobilized nitroxide spins and 2 ``dilute'' mobile nitroxide spins is roughly $2.2\times10^{13}$ each. .........}

\bigbreak

\np{We then measure the in-operando cwEPR from 0.94~V towards 0.3~V thereby reducing TEMPO\textsuperscript{$+$} $\rightarrow$ TEMPO\textsuperscript{$\bullet$}. Here we see that reduction occurs over a larger potential range (as compared to oxidation) c.f. Figure \ref{fig:Figure_4}b. This could be due to a changing morphology of the p-DiTS films which decreases electrolyte access to the TEMPO groups in the film matrix.}

\np{While the EPR signature of different polymeric NiSalen complexes has been reported by Dmitrieva \textit{et al} \cite{Dmitrieva2018}, we do not observe the EPR signature of the NiSalen backbone in our in-operando study of p-DiTS. Dmitrieva \textit{et al} also noted that under high doping levels of the polymer the charges could aggregate to form diamagnetic states and under high potential polymers could be doped with two electrons per monomer also forming dications which are also diamagnetic. Another plausible explanation for the absence of a NiSalen EPR signals, could simply be low intensity signals masked by the much more intense nitroxide signals. However, our estimate of the number of spins per monomer units in this film is 2.04, suggesting only signals corresponding to nitroxide radicals where there are two unpaired electron per monomer unit.}

%\ik{
%Voltage-dependent cwEPR signals were observed upon oxidation and reduction of a p-DiTS film inside the tube-based electrochemical cell (Figure \ref{{fig:Figure_4}}). The initially reduced cell was oxidized by applying potentials from 0.30$\,$V to 0.94$\,$V (vs. Ag) that is shown in the "oxidation" part of the plot a). The highest applied potential was 0.94$\,$V that is shown in a dashed line in the middle of the plot. Then the cell was reduced back to 0.30$\,$V (that part is labeled as "reduction"). Initially, the p-DiTS film has a high density of TEMPO$^{\bullet}$ radicals that introduces dipolar coupling between the neighboring spins in the film which broadens the cwEPR line. During the oxidation, the intensity and the shape of the spectral line is stable until 0.75$\,$V (slice I). Then, between 0.78 and 0.90$\,$V, the cwEPR signal decreases steeply.\\ 
%}

%\bigbreak
%
%\ik{
%A quantitative analysis of the spectra is shown for the corresponding potentials on subplot b). The background spin concentration for the EPR spin count was determined from noise analysis described in the SI. The total number of the electrochemically active charges was extracted from the CV of the cell and is plotted in b) for comparison. The EPR signal was averaged within the oxidized state of the cell (region II, 0.90$\,$V to 0.94$\,$V to 0.82$\,$V). The averaged signal of the oxidized cell, scaled up by a factor of 10 is shown in slice II. It was decomposed into two components with numerical simulations that is described further in the text. The reduction of the cell appeared with a delay in applied potentials, that is seen in the asymmetry of the region II with respect to the highest applied potential. The cwEPR signal reappears at 0.78$\,$V, which is 120$\,$mV lower than the oxidation steep at 0.90$\,$mV. The reduction of the cell leads to a complete recovery of the signal (slice III) that is seen in the recovered spin count.\\ 
%}

%\bigbreak
%\ik{
%The number of spins is changing from 6.43e+14 in the reduced state to 5.40e+13 in the oxidized state (with the 20\% uncertainty). The number of electrochemically active TEMPO groups determined from the cyclic voltammogram is 7.6e+14. The background spin count is 6.19e+12. The amount of substance in the film was found from the charge of the film with the Faraday electrochemical constant. The charge of the film was found by integrating the cathodic branch of the cyclic voltammogram of the cell vs. time. The mass of the film was found to be 1.6$\,\mu$g using the calculated amount of substance and the molar mass of DiTS of 863 g$\,$mol$^{-1}$. The electrochemistry measurements of the mass of the film combined with the spin count data yield a spin density of 2.04 spins per DiTS monomer, that correlates to the two charge bearing groups per DiTS monomer and indicates that the charge of the film is composed solely by the charge carrying nitroxide groups. 
%For the measurements shown in Figure~\ref{{fig:Figure_4}}, the p-DiTS film was rinsed after the polymerization. That had removed most of the non-polymerized fragments from the film as almost no free TEMPO signal is observed in region II.\\
%}
%\bigbreak
%
%The change of the oxidation potential of the film together with the change of the number of electrons determined from quantitative cwEPR measurements allows one to measure the electric energy density of the charged film. In the described experiment, the number of electrons varies by 5.90e+14 and the electric potential changes between 0.30 and 0.94$\,$V, that yields the energy density of the film of 5.2E-3 W~h~g$^{-1}$.

\bigbreak

\subsection*{cwEPR -- a Tool to Determine the Quality of ORB Polymer Films}

\np{
	One of the most important parameters for designing any battery system is to check the purity of the active layers (cathode and anode). As such cwEPR can play a vital role in investigating the purity of the active polymer layer in ORBs by distinguishing between molecules which are redox active and contributing to the capacity and those that are not. Furthermore, cwEPR can be used to understand the nature of the contributing and non-contributing molecules by changes in spectral signature, based upon which improvements can be made to the battery or active layer materials.
}

%CwEPR can be utilized for investigating the availability of the electrochemically-active molecules in the active polymer layer in ORBs. Furthermore, cwEPR can be used to understand the nature of the contributing and non-contributing molecules by changes in the spectral line shapes.

\bigbreak

\subsubsection*{Release of TEMPO fragments into Electrolyte}

\np{During the electropolymerisation/deposition process it is possible that some monomeric DiTS which has not polymerised is trapped in the pores or interlayer structure of the growing polymer film. These monomers are then perhaps not redox active as they are not directly linked via a long conductive backbone but more crucially these monomers can be released into the electrolyte where they most definitely do not contribute to the battery performance. To investigate this phenomenon we measure the in-operando SEC cwEPR of a 150~CV cycle p-DiTS film which did not go through the rinsing procedure.}

\bigbreak

\np{	
While electrochemical quartz crystal microbalance measurements can hint towards loosely bound oligomers or trapped monomers being released from a newly grown polymer film after rinsing or initial CV cycling, it can not distinguish between 1) oligomers and monomers released or 2) as relevant in our case between mono-TEMPO fragments and di-TEMPO monomer units. cwEPR can give information in both of the mentioned cases, for 1) though changes in rotational correlation time ($\tau$) between a monomers (faster $\tau$) and oligomers (slower $\tau$) due to a increasing molecular size and for 2) spectral changes from three-lines to five-lines due to exchange coupling which is present in the di-TEMPO monomer but not mono-TEMPO fragments. 
}

\bigbreak
%During the electropolymerisation/deposition process, it is possible that some monomeric DiTS which has not polymerised, is trapped in the pores of the growing polymer film. These monomers are then perhaps not redox active as they are not directly linked to the electrode via a long conductive backbone. The monomeric DiTS can be released into the electrolyte where they do not contribute to the battery performance. The release of the monomer fragments can be detected by cwEPR, as in the reduced state they are paramagnetic and exhibit a characteristic three-line cwEPR spectra. To detect the released nitroxide fragments, we measured the in-operando SEC cwEPR of a 150~CV cycle p-DiTS film which did not go through the rinsing procedure and contained significant amount of the not polymerized fragments. 

%\bigbreak

%In the \textit{in-operando} SEC cwEPR (\textbf{Figure \ref{fig:Figure_5}}a) of a raw p-DiTS film, at 0.3~V there is a single derivative corresponding to the exchange/dipole broadened densely packed film. As we oxidised the film by applying 0.9~V the signal intensity decreases but by a relatively small amount, however at 1.2~V the broad component vanishes which corresponds to a full oxidation of the densely packed film. The residual signal in the 1.2~V cwEPR spectrum is a 3 line spectrum which indicates that there are mono-TEMPO fragments which have dissolved into solution. 
\np{
In our in-operando SEC cwEPR (Figure \ref{fig:Figure_5}a) of raw p-DiTS film we observe that at 0.3~V we see a single derivative corresponding to the exchange/dipole broadened densely packed film. As we oxidized the film by applying 0.9~V the signal intensity decreases but by a relatively small amount, however at 1.2~V we see a full oxidation of this densely packed film. The residual signal in the 1.2~V cwEPR spectrum is a 3 line spectrum which indicates that there are mono-TEMPO fragments which have dissolved into solution. The presence of mono-TEMPO in electrolyte (also confirmed by measuring the used electrolyte) suggests degradation of our DiTS molecules. This degradation and release of nitroxide fragments into electrolyte would be a significant problem for ORBs and as can be seen from the reduction steps from 1.2~V to 0.9~V and 0.3~V, the one line spectrum recovers but not to the same initial intensity and the three-line component is seen superimposed. Further investigation of this three-line feature could lead to information about the size of the fragments by comparison to TEMPOL molecules in electrolyte, the mono-TEMPO succinyl moiety and mono-TEMPO NiSalen in electrolyte solution and studying their rotational correlation times. The exact size and nature of the nitroxide fragments is beyond the scope of the current study, but can be identified with EPR.
}

\bigbreak

\begin{figure}[h!]
	\includegraphics[width=0.5\textwidth]{Figure_5}
	\caption{\textit{in-operando} SEC cwEPR of two different p-DiTS films a) 50~CV cycles, (Top) showing the appearance of a superimposed three-line feature, corresponding to mobile nitroxide groups, with increasing oxidation potential. The three-line feature is then electrochemically inactive and remains in during the reduction cycle. The rest electrolyte (Bottom) shows the same three-line feature suggesting release of nitroxide fragments into the electrolyte. b) 200~CV cycles, (Top) showing the appearance of a superimposed three-line feature which is electrochemically active and after oxidation does not re-appear in the reduction cycle, suggesting surface oligomers which are initially mobile but later polymerise to the main film or being released to the electrolyte in the EPR-silent, oxiized state. (Bottom) rest electrolyte shows negligible signal of nitroxide spins.}
	\label{fig:Figure_5}
\end{figure}

\bigbreak

\subsubsection*{Re-Polymerization and missing spins}

In a 500~nm-thick p-DiTS film which was made using 200~CV cycles and underwent the cleaning procedure, the cwEPR signal of the fully reduced state (0.3~V) shows a significant increase in the signal intensity due the increased thickness. Upon applying an oxidation potential of 0.9~V the one-line feature decreases as it is expected from the oxidation of the p-DiTS film. At the same time, the satellite peaks emerge that correspond to mobile nitroxides fragments being released into the electrolyte. After further oxidation to 1.1~V, the one-line feature has disappeared which corresponds to the full oxidation of the film. A negligible three-line spectrum is still present, though much less than can be observed in the 0.9~V oxidation spectrum.\\

\bigbreak

\np{
Reducing the film back from 1.1~V to 0.9~V we see the partial recovery of the signal and similarly reducing to 0.3~V still only results in partial reduction of the film/ recovery of the cwEPR signal. Since it is well known that TEMPO exhibits a reversible redox reaction from the TEMPO$^{\bullet}$ $\rightarrow$ TEMPO$^{+}$ $\rightarrow$ TEMPO$^{\bullet}$, it is unlikely that the decrease in overall signal intensity in the Figure~\ref{fig:Figure_3}~b during the \textit{in-operando} SEC cwEPR is due to a chemical reaction. Rather the more plausible explanation is that oligomers in this thick film are oxidised and then detach from the electrode thereby no longer being electrically accessible for reduction and also not appearing in the only electrolyte spectrum as TEMPO$^{+}$ is EPR silent. Furthermore the presence of satellite peaks corresponding to mobile nitroxides, is different to the case seen in Figure~\ref{fig:Figure_3}~a. The presence of a superimposed three-line spectrum (Figure~\ref{fig:Figure_3}~b) 0.9~V oxidation) which can be electrochemically oxidised (Figure~\ref{fig:Figure_3}~b) 1.1~V oxidation), but does not reappear in the 0.9~V reduction spectrum, suggest that this signal correspond to oligomers which are electrically connected to the electrode but are initially mobile. Such behaviour could be expected of oligomers near the film surface, where one end can be mobile and the other end connected through the main film structure to the electrodes. After initial oxidation up to 1.1~V the mobile end of the oligomer goes through a polymerisation reaction where by it is more rigidly connected to the main film structure, thereby removing the three-line mobile nitroxide signal from the cwEPR.
}

\bigbreak

The spectral deconvolution into the separate components and the evolution of the spectral line width is presented in the SI, \textbf{Figure~\ref{fig:S8_SI}}

\bigbreak
\np{SEC cwEPR can therefore be a powerful tool to understand differences in film behavior with different thicknesses and changes in the local film structure. The above study shows three different yet important pieces of information that SEC cwEPR can yield which is not readily available from conventional electrochemistry techniques. Firstly, though SEC cwEPR we can not only see/quantify that some active electrode material is lost to the electrolyte but we can also see if it is in a oxidized or reduced state. Secondly, we can observe through clear spectral signatures nitroxide pendent groups at the surface which are mobile yet still electrochemically active and thirdly we can observe how these initially mobile groups subsequently become immobilized as the oligomer backbone (re-)polymerizes to the main film structure. Such subtle difference in film structure and behavior during operation and vital for understanding and improving organic materials for active electrode materials in ORBs and would provide vital information for synthetic chemistry and electrochemistry but the ORB field in general.
}

\bigbreak

\subsection{Degradation of p-DiTS upon Charge-Discharge Cycling}\label{Degradation_study}

\begin{figure}
	\includegraphics[width=0.95\textwidth]{Figure_7}
	\caption{a) Evolution and component decomposition of cwEPR spectra for a tube-based, 150~CV p-DiTS ORB upon 36 charge-discharge cycles. b) Quantitative analysis of the separated spectral components. Quantitative analysis of the CV curves. c) Cyclic voltammograms of the cell, recorded with respect to Ag/AgCl RE at a rate of 5 mV/s inbetween the EPR scans, in the modified tube inside the ER~4122-SHQE microwave resonator.}
	\label{fig:Figure_7}
\end{figure}

\ik{
The p-DiTS based ORBs were shown to lose their capacity under reprated charge-discharge cycling, but the mechanism of degradation is not yet understood on the molecular level \cite{vereshchagin2020_battandsupercaps}. We performed a series of charge-discharge cycling on a tube-based p-DiTS electrochemical cell and recorded CV curves together with the cwEPR spectra during the cycling.\\
}

\bigbreak

\ik{
For the degradation experiment, a 389$\pm$40~nm thick p-DiTS film (see Figure~\ref{fig:S_SAMPE_INFO}) was deposited onto the on-substrate WE of the tube-based cell. The cell was inserted into the cwEPR spectrometer, at $H\textsubscript{res. center}$~~12~mm where it underwent a series of 36 charge-discharge cycles at a rate of 5$\,$mV/s. The CV curves were recorded for each cycle. The cwEPR spectra were recorded after each fourth cycle for the fully reduced cell, that is at the left-most point on the corresponding CV.\\
}

\bigbreak

\ik{
\textbf{Figure~\ref{fig:Figure_7}} represents the cwEPR spectra (a) and the CV curves (c) during the cycling. The cycles inbetween are shown in Figure~\ref{fig:S_Deg}. During the cycling, the CV curves decrease in intensity that indicates the reduction in the number of electrochemically active charges in the film. The number of electrochemically active charges was determined for each cycle by integrating the reduction branch of the corresponding CV, and plotted in Figure~\ref{fig:Figure_7}(b) as "CV".\\
}

\bigbreak

\ik{
The cwEPR spectra in Figure~\ref{fig:Figure_7}(a) demonstrate initially a release of the dilute component and a decrease in intensity upon further cycling. 
%The deconvolution of the spectra into the "packed" and "dilute" components allows for distinguishing between the overlapping EPR signals that come from DiTS molecules situating in two distinct environments (cf. Section~\ref{subs:classification}). 
The three-line, "dilute" component corresponds
% to the reduced monomeric DiTS or (more general) 
to the TEMPO$^{\bullet}$ fragments that detach from the film in the reduced state and get released to the electrolyte. The "dilute" component is emerging only during the first eight cycles. Meanwhile, the amplitude of the CV in (c) stays as high as 100~$\mu$A during the first eight cycles. That suggests that the initial release of the reduced, EPR active TEMPO fragments does not affect the capacity of the cell and can be excluded as a possible degradation pathway.\\ 
The broad, single-line "packed" component, on the contrary, does decrease for the later cycles, together with the CV. Therefore, the degradation of the electric capacity of the polymer film is connected with the decrease in the "packed" component and is not accompanied by the release of any EPR-active products to the electrolyte.\\
}

\bigbreak

\ik{
We performed a quantitative analysis of the cwEPR spectral components (cwEPR spin count) and obtained the number of paramagnetic species in the reduced film at each degradation step. We also calculated the number of electrochemically active electrons by integrating the corresponding CV (CV electron count). The cwEPR spin count and the CV electron count are shown in Figure~\ref{fig:Figure_7} b). 
	%The relative contribution of the "dilute" component to the spin count is neglegibally small.\\
	For the fresh film, "0 CV", CV electron count coincides with the cwEPR spin count within the measurement uncertainty. However, by the last, 36$^{th}$ cycle, the CV electron count is significantly smaller than the cwEPR spin count.
}

\bigbreak

\ik{
All cwEPR spectra in Figure~\ref{fig:Figure_7}a) were recorded when the film was brought to the reduced state. There, all electrochemically active TEMPO$^\bullet$ fragments within the film are also EPR active. Additionally, all electrochemically inactive fragments that happen to be in the reduced, TEMPO$^\bullet$ state, also contribute to the cwEPR signal. Thus, neglecting the insignificant amount of the released monomer fragments, the cwEPR signal at each charge-discharge cycle contains two contributions: 1) the densely packed TEMPO$^\bullet$ fragments that are electrochemically active and were brought to the reduced state by discharging the cell, and 2) the densely packed TEMPO$^\bullet$ fragments that are not electrochemically active and always stay in the TEMPO$^\bullet$ state, regardless of the applied potential.
% That, together with the study of the molecular backbone, NiMeOSalen, in Section \ref{Section:backbone} and the operando spin count in Firure \ref{{fig:Figure_4}} suggest that the EPR signal from a p-DiTS film is comprised only by the nitroxide fragments inside the film. 
Thus, the overall spin count on the cwEPR spectra in Figure~\ref{fig:Figure_7}b) represents the number of TEMPO$^\bullet$ fragments inside the reduced film, even though some fraction of them is not electrochemically active and some insignificant fraction of them had detached from the film.\\
}

\bigbreak

\ik{
The "packed" spectral component decreases with cycling. That indicates, some paramagnetic parts of the film (TEMPO$^{\bullet}$) become diamagnetic (TEMPO$^{+}$) and stay diamagnetic when the cell is brought back to the reduced state. These EPR silent fragments may either stay in the film or to be released into the electrolyte. 
From the EPR only perspective, it is impossible to determine whether the oxidised fragments stay in the film or leave it. However, in the previous study, no loss of mass of a p-DiTS film was observed under a microbalance-monitored charge-discarge cycling \cite{vereshchagin2020_battandsupercaps}. This excludes a scenario where the film would lose its fragments upon repeated cycling. The diamagnetic species in the electrolyte may be brought to the paramagnetic state with electrochemical reduction, and then detected with cwEPR.\\
}
\ik{
A more rapid decrease of the CV electron count indicates that some regions of the film become electrochemically incative upon cycling, while staying cwEPR active. That is, the  TEMPO$^{\bullet}$ fragments lose the electrical connection to the rest of the film and stay in the reduced state. This allows one to assume that the reduced fragments stay in the film, but lose the electrical connection to the electrode.\\
}

\ik{
Combining the dynamics of the cwEPR spectra and the CV, we propose that the degradation of the electrochemical properties of the film is caused by a fragmentation process that leads to the formation of electricaly isolated domains within the film. Both charged and discharged domains may form inside a p-DiTS cathode film upon its repeated charging and discharging. In other words, the p-DiTS based ORB loses its capacity in both charged and discharged states, due to the gradual loss of electrical connection upon cycling that leads to a formation of electrically isolated molecular domains. The faster degradation of the CV as compared to the cwEPR indicates that more domains are formed in the reduced, TEMPO$^{\bullet}$ state than in the oxidised, TEMPO$^{+}$ state.\\
}

\bigbreak

\subsubsection*{Ex-Situ SEC cwEPR on p-DiTS at 150~K}
\label{Section:backbone}
While in-operando or even in-situ SEC cwEPR provides information about redox systems under the natural operating environments, these measurements for the batteries are limited by the time for which potentials can be applied before degradation of the active materials is observed. For the p-DiTS cells, \textit{in-operando} SEC cwEPR measurements were carried out at room temperature while holding the redox potential of interest for $\sim$200~seconds and accumulating the cwEPR spectra. The S/N ratio for this setup is limited by the time of scan and by the low population difference of the spin levels due the high temperature of the sample. We therefore perform \textit{ex-situ} measurements on a 50~CV cycle film at cryogenic temperatures, in this case at 150~K. Those \textit{ex-situ} measurements involved bringing the cell to the chosen redox potential in a beaker, outside the spectrometer, using the setup shown in Figure~\ref{fig:Figure_1}~c. The applied potential was held for 100--200~seconds, then the electrodes were disconnected from the cell. After that, the substrate with the on-substrate electrodes was sealed in an ordinary, 5~mm OD quartz tube and flash-frozen in LN$_2$. We find that this method can hold the cell at different redox state regions but not necessarily the exact redox potentials as the charge seems to equilibrate to different overall states. Nevertheless, the aim of the \textit{ex-situ} measurement was not to see a step by step oxidation and reduction, but to reach representative redox states which were also seen in the \textit{in-operando} measurements (Figure~\ref{fig:Figure_4}), i.e. fully reduced, partially oxidised and fully oxidised states, with a higher S/N ratio.\\

\begin{figure}[h!]
	\includegraphics[width=1\textwidth]{Figure_6}
	\caption{\textit{Ex-situ} oxidation series at 150~K on a) p-DiTS (black), showing a one line derivative at 0.1~V which with increasing potential decreases in intensity as nitroxides groups are reduced and changes in spectral shape towards a dilute immobilised nitroxide signal. Simulations (green) show there two components contributing to the signal at each potential, a potential dependent one line spectrum (red) and a immobilised nitroxide spectrum (blue) which is potential independent. b) p-NiMeOSalen revealing the backbone paramagnetic state which has rhombic $g$-tensor and shows small spectral changes as a function of potential.}
	\label{fig:Figure_6}
\end{figure}

\bigbreak

We therefore measure the \textit{ex-situ} SEC cwEPR at different oxidation potentials three of which are shown in \textbf{Figure~\ref{fig:Figure_6}}a. At 0.1~V the p-DiTS film is in the reduced state and exhibits a densely packed one-line feature, at 0.6~V and 0.9~V the overall signal decreases significantly and a similarly broad yet distinctly different signal appears, that was singled out with the numerical simulations of the spectra. The signal at each of the three potentials has contributions from the nitroxides/TEMPO groups in two distinctly different environments. The broad one-line feature coming from densely "packed" dipole and exchange coupled nitroxide spins, and "dilute" immobilised nitroxides which have a rhombic $g$-tensor. We can simulate the measured spectra at each oxidation potential using EasySpin \cite{Stoll2006}, using a combination of a packed and dilute contribution with changing relative weights. We find that with increasing oxidation potential the one-line packed nitroxide feature decreases but the intensity of the dilute immobilised nitroxide does not change. The film in this ex-situ study did not contain any electrolyte, so both spectral components correspond to the nitroxide fragments situated in the film. That suggests that the p-DiTS films have some nitroxide spins which are not electrochemically active yet are still a part of the film.

\bigbreak

\subsubsection*{\textit{Ex-Situ} SEC cwEPR on p-NiMeOSalen at 150~K}

The conductive backbone of p-DiTS i.e. NiSalen is also a redox active material and has been shown to give EPR active states at different oxidation potentials (in-situ). However, in the p-DiTS system we did not detect the paramagnetic state of the backbone. In order to observe the paramagnetic states of the backbone, we measured the \textit{ex-situ} SEC cwEPR at 150~K of p-NiMeOSalen which closely resembles the core backbone structure of p-DiTS but does not contain any TEMPO pendent groups. Dmitrieva \textit{et al.}~\cite{Dmitrieva2018} have studied Nickel Salen polymers using \textit{in-situ} SEC cwEPR, among which was NiMeOSalen. Dmitrieva \textit{et al.} observed three different paramagnetic species in NiMeOSalen, depending on the oxidation potential, with differing $g$-tensor/factor values. We only investigated p-NiMeOSalen \textit{ex-situ} and saw a change in the $g$ tensor, but no significant change in the signal intensity with respect to the applied potential. 

\bigbreak

The SEC cwEPR of p-NiMeOSalen at 150~K at a few key potentials is shown in \ref{fig:Figure_6}, while the full dataset and the measured components of the g tensor with respect to the applied potential are shown in \ref{fig:S4}. The 0~mV cwEPR spectrum correponds to $g$=[2.0235 2.0094 1.9984] that is in agreement with the values measured by Dmitrieva \textit{et al.} ($g$=[2.0240 2.0094 1.9990]), considering that our measurement was done \textit{ex-situ}.\\
The further evolution of the cwEPR spectra does not correlate with the applied potential and no trend is seen for the components of the $g$ tensor. That justifies that the backbone can not store the electrical charge.\\
Our in-situ SEC cwEPR measurements on p-NiMeOSalen films at room temperature did not yield appropriate S/N ratio.\\

\bigbreak

\subsection{Pulse EPR on p-DiTS - a step closer to SEC pEPR}

\bigbreak

\np{To the best of our knowledge SEC pEPR has so far been illusive, primarily due to fast spin relaxation but also microwave damping from the presence of electrodes and high dielectric electrolyte, at room temperature. As such the advanced techniques of potential dependent pulse EPR have so far not been realised. Potential dependent, or SEC pEPR can yield a wealth of information not available from SEC cwEPR. Two important pEPR methods useful for any electrochemically active material is PELDOR and Hyperfine spectroscopy such as ENDOR or HYSCORE. Analysis of PELDOR traces give information regarding electron-electron distances, while hyperfine spectroscopy gives information regarding the electron spin (de-)localisation on the molecule. The mentioned methods could give insight into degradation mechanisms and environment of the electrochemically inactive spins in pDiTS but more generally ORB materials. Realisation of SEC pEPR, particularly on film electrodes, would also make these techniques available for other research field such as the study of hydrogenases and enzyme or metal catalysis to name a few.}

\bigbreak

\np{As a first step, we perform pEPR on p-DiTS measured at 20~K in its fully reduced state when deposited on both the WE and CE. We use the Hahn Echo sequence ($\frac{\pi}{2} - \tau - \pi - \tau - Echo$) to record the integrated echo intensity as a function of the swept static magnetic field. The obtained echo detected field swept (EDFS) spectrum is shown in .... Here we see the typical immobilised uncoupled/weakly coupled nitroxide powder spectrum, where the spectral width at X-band frequencies is given by $2A\textsubscript{zz}$ where $A\textsubscript{zz}$ is the Z component of the hyperfine tensor.}

\bigbreak

\np{The observation of non-coupled “dilute” nitroxide spins in the EDFS spectrum and not the “packed” coupled spins is due to the slower relaxation times of non-coupled spins compared to coupled spins. The EDFS spectrum in .... is likely a result of the same electrochemically inactive spins observed in the ex-situ pDiTS redox series (Figure \ref{fig:Figure_6}). Potential dependent SEC pEPR measurements are ongoing to confirm this hypothesis and furthermore PELDOR and hyperfine spectroscopy could lead to critical information about the local environment which causes these spins to be electrochemically inactive.}

\bigbreak

\np{Further measurements are also ongoing to detect the pEPR spectrum corresponding to the strongly coupled “packed” nitroxide spins where potential dependent SEC pEPR will have the most impact. These SEC pEPR measurements on the “dilute” and “packed” spins will form the basis of a future publication.}


\section{Conclusion and Perspectives}
Our results show that electron spins can be used as unique local probes for elucidating redox reactions associated with charging and discharging of ORBs with active electrode materials prepared by electropolymerisation. EPR spectroscopy can thus be employed to analyse these processes, identify performance-limiting loss mechanisms, and eventually help developing strategies for making polymer batteries powerful contenders on the path towards sustainable electrochemical power sources. 


\bigbreak

The fact that we can observe a strong influence of the electrical potential on the in-situ EPR spectrum at ambient conditions provides the basis for systematically studying the redox-active sites in different states of the active electrode materials. For instance, changes in the film structure, e.g.\ caused by degradation, can be detected in in-situ cwEPR experiments (cf.\ Figure~\ref{fig:Figure_5}) as nitroxides in different environments are associated with unique and characteristic EPR signatures (cf.\ Figure~\ref{fig:Figure_2}). Much more detailed and specific information about degradation processes can, however, be extracted from advanced (pulse) EPR experiments which usually require low temperatures, mostly to increase spin-relaxation times. We believe that the possibility of performing low-temperature pEPR measurements is the key factor that is mandatory for a widespread use of advanced EPR techniques for ORB research.\\


\bigbreak


The electrochemical cells for in-situ studies, as described in this article, were be designed in such a way that they fit into a standard EPR tube. They are thus compatible with (low-temperature) experiments using conventional pEPR instrumentation. This opens up the intriguing perspective of exploiting the coupling between neighbouring electron spins as well as the hyperfine coupling between electron and nuclear spins in close vicinity to determine the coupling (and thus the distance) between redox-active centres and the degree of delocalisation for charge carriers on the conjugated polymer backbone. In particular, it is conceivable to find out whether inactive redox centres, which are not changing their redox state upon charging and discharging, can be found in clusters or are rather isotropically distributed in the cathode films. This provides the possibility to identify inactive redox sites that reduce the maximum attainable cell capacity. The pEPR measurements can either be performed ex-situ (without electrolyte, cf.\ Figure~\ref{fig:Figure_6}) or on samples that are flash-frozen while the desired electrical potential is applied to the cell.\\

\bigbreak

The experiments presented here are by no means limited to the electrochemical cells used in this study. Specifically, they could also be used for investigating all-organic solid-state batteries such as the recently reported and very promising cells comprising polymeric anode and cathode materials as well as gel electrolytes based on ionic liquids \cite{muench2021_jps}.\\

\bigbreak

A strategy for overcoming sensitivity limitations of advanced EPR experiments is to combine them with electrical readout schemes and make Electrically Detected Magnetic Resonance (EDMR) applicable to study redox and transport processes in polymer batteries. EDMR possesses a sensitivity many orders of magnitude higher than conventional EPR measurements.\cite{mccamey2006_apl} The possibility of using EDMR techniques for battery characterisation at ambient conditions would also open up entirely new opportunities for imaging techniques. EDMR imaging could provide a detailed picture of the spatial distribution of active and inactive redox centres and, in general, usefully complements information obtained through imaging techniques based on conventional EPR~\cite{sathiya2015_natcomm} or nuclear magnetic resonance.\cite{ilott2016_pnas}


\section{Experimental Methods}
The cwEPR spectra were recorded at X-band frequencies ($\sim$9.4~GHz) using a laboratory-built EPR spectrometer. The magnetic field was regulated using a Bruker BH15 field controller and monitored with a Bruker ER 035M NMR gaussmeter while the Bruker ER048R microwave bridge was used for microwave generation (200~mW source) and detection (diode-detection). The static magnetic field was modulated at 100~kHz and lock-in detection was carried out using a Stanford Research SR810 lock-in amplifier in combination with a Wangine WPA-120 audio amplifier. The lock-in detection leads to the derivative spectra. $g$-factor calibration was additionally done using a N@C\textsubscript{60} powder sample with $g=$2.00204. SEC cwEPR spectra were taken using the setup described earlier, using a on-substrate electrode setup, with a polymer film deposited on the WE using electropolymerization, and modified EPR tube where the electrolyte was 0.1~M solution of Et\textsubscript{4}NBF\textsubscript{4} in ACN ($\epsilon\approx40$, T\textsubscript{melt}$=$227~K). Control over the redox potential of the cell was achieved by using a Keithley 2450-EC potentiostat.\\

\bigbreak

CV curves were measured either in a beaker using the on-substrate WE, a coiled Pt CE and an Ag/AgNO\textsubscript{3} RE with the above mentioned electrolyte (Figure \ref{fig:Figure_1}c) or in a modified EPR tube for in-operando measurements (Figure~\ref{fig:Figure_1}d), with the on-substrate WE and CE and a Ag or Ag/AgCl wire as a RE with the same electrolyte. In-operando SEC cwEPR measurements were done by holding the sample at the redox potential of interest for 200~s using the chronoamperometry mode on the Keithley 2450-EC potentiostat, while the cwEPR scans were accumulated. The 200~s scans allowed for cwEPR spectra with a suitable S/N ratio without causing any noticeable degradation to the electrochemical cell up to a potential of 950~mV.\\

\bigbreak

The potential-dependent cwEPR spectra were recorded at room temperature for in-operando measurements and at 150~K for ex-situ measurements without electrolyte. The low temperature measurements provide an increased S/N ratio for the cwEPR spectra and therefore allowed for a better decomposition of the spectral components contributing to the signal. The potential-dependent in-operando cwEPR spectra were measured for p-DiTS, while the ex-situ cwEPR spectra were measured for both p-DiTS and p-NiMeOSalen films.

\np{Details of all solvents, salts and structure and compositions of all materials must be included!! So salt, ACN, PC, silver wire, platinum wire, metals used for thermal evaporation etc for sure we need the details on the bottle, Isopropyl alcohol and Ethanol also maybe.}

\bigbreak

\medskip

\textbf{Supporting Information} \par %Please delete the Suppporting Information statement if it is not applicable. Please supply Supporting Information in another file. Supporting information should not be provided in .tex format
Supporting Information is available from the Wiley Online Library or from the author.



% Acknowledgements
\medskip
\textbf{Acknowledgments} \par %delete if not applicable))
Please insert your acknowledgments here

% References
\medskip

% Use the following code if you wish to generate your bibliography with BibTeX;
% replace the string "MSP-template" below with the name(s) of
% the BibTeX data base(s) you want to use.
% The resulting bibliography-output (the content of the .bbl file)
% must be pasted back into this file before submission.
% Please also include your BibTeX data base file(s) in your submission
% so that we can re-run BibTeX if necessary.
%
%\bibliographystyle{MSP}
%\bibliography{MSP-template}


\bibliographystyle{MSP}
\bibliography{ORB}% Produces the bibliography via BibTeX.






% Figures/tables and captions
% Permission statements are required for all figures reproduced or adapted from previously published articles/sources. Please also ensure that all necessary permissions to reproduce images have been received
% Please remove these statements for original figures


% Please provide Biographies and photos for Essays, Feature Articles, Progress Reports, Reviews, and Perspectives for those authors who should be highlighted  
% These should be at most 100 words long
% For other article types this section can be removed
% Photographs should be 40mm broad and 50 mm high

\newpage

\section{Supporting Information}
% This resets the figure counter to 0 and adds a 'S' in front of the numbers. Compile twice to get it correct in the references
\renewcommand{\thefigure}{S\arabic{figure}}
\setcounter{figure}{0}

\subsection{Fabrication of the Modified Tube}
\par
The volume of the cell had to be minimised, as the high dielectric constant of the electrolyte does not allow for the critical coupling of the resonator. The volume of the electrolyte was minimized by flattening the tip of the quartz (melting point $T_m=1660-1710\,^{\circ}$C) tube accomodating the cell. For that, a tungsten (W, $T_m=3422\,^{\circ}$C) rod of a 3.70$\times$1.35$\,$mm rectangular cross section, narrowing to the bottom of the tube at a wedge angle of $\alpha \approx 5 ^{\circ}$, was inserted into the tube (Figure \ref{fig:S_ROD}). The tube was evacuated and filled with He to a pressure of 100 mbar. The tip of the evacuated tube was molten around the thermally expanded rod with a Hydrogen-Oxygen burner ($T_f = 3080\,^{\circ}$C). The modified tube was then connected to atmosphere and the cooled, contracted rod was removed and later reused for modifying other tubes. The tungsten rod was fabricated in the student workshop of the department of physics of the Free University of Berlin from an electrode of a used flash lamp of the Quanta-Ray Pulsed Nd:YAG laser.\\

\begin{figure}[h!]
\includegraphics[width=0.95\textwidth]{Figure_S1}
\caption{Modifying the shape of the fube with a W forming rod.}
\label{fig:S_ROD}
\end{figure}

\par The modification of the tip of the tube provided a decrease in the volume of the electrolyte downto 45$\,\mu$l, while the cylindrical part of the tube allowed for fitting the substrate connector and the wiring in. That allowed for the room temperature EPR mesurements on the electrochemical cell during charging and discharging. Extensive care had to be taken to reproduce the cv curve in the modified tube as the limited volume of the electrolyte was distributed over the large surface area.\\

\bigbreak

\subsubsection{Ag/AgCl Reference Electrode}
\label{Ref_electrode}
The Ag/AgCl reference (RE) electrode was formed by coating a layer of AgCl onto a $\varnothing\,$250$\,\mu$m Ag wire galvanically from a 1$\,$M aqeous solution of KCl. For that, the Ag wire and a Pt wire were connected as a cathode and an anode respectively to 10$\,$ml of the KCL solution and a DC voltage of 500$\,$mV was applied for 7500$\,$s which resulted in a dark AgCl layer deposited onto the Ag cathode. In the beginning of the galvanic deposition, the current through the KCl solution was changing in a jumpy manner, that indicated the growth of the AgCl layer. The galvanic process was stopped after the current was stable for at least 300$\,$s.\\



\newpage

\subsection{Q factors and Signal to Noise ratio for the Modified Tube}\label{Q-factor_SNR}
The electrochemical cells based on the modified tube allow for higher EPR signals as compared to the normal X-band EPR tube. Figure \ref{fig:S_SAMPE_INFO} represents two sets of measurements that describe the enhancement of the EPR signal for the modified tube. When the cell is inserted closer to the resonator's center, the Q factor of the resonator lowers and the EPR signal increases, because the intensity of the magnetic component of the microwave field B$_1$ is higher in the center of the resonator. When the cells were inserted closer to the center of the resonator, the Q factor went below $\approx$~600, so that the signal arm of the spectrometer could not be critically coupled and no EPR measurements were possible. Depending on the volume of the electrolyte in the cell, it can be inserted more or less deep into the resonator.\\
\par The modified tube could be inserted 7 mm deeper into the microwave resonator (see Figure \ref{fig:S5}). The EPR signal in the modified tube is increased as compared to the normal tube at the same height.\\

\begin{figure}[h!]
\includegraphics[width=0.95\textwidth]{Figure_S2}
\caption{cwEPR signal intensity and Q factors for conventional EPR tube and the modified tube at different hights from the center of the Bruker SHQ resonator.}
\label{fig:S5}
\end{figure}


\newpage
\subsection{Sample Information}

\begin{figure}[h!]
\includegraphics[width=0.67\textwidth]{Figure_S3}
\caption{Sample information. Deposition curves, cycling curves, EPR parameters, figure references.}
\label{fig:S_SAMPE_INFO}
\end{figure}

\subsection{Spin Count Data}

\begin{figure}[h!]
\includegraphics[width=0.88\textwidth]{Figure_S4}
\caption{Spin count during the (secondary) release of the monomer from a thick (150 deposition cycles), untreated p-DiTS cell. Before this measurement, the cell has already exhibited the (primary) release of the monomer during the initial charge-discharge cycle, then it was rinsed in Acetonitrile and the electrolyte was exchanged. The three EPR lines in the oxidized cell correspond to the released nitroxide fragments that freely tumble in the electrolyte. a): CwEPR spectra during the charge-discharge cycle of the untreated p-DiTS cell that releases the monomers. b): number of unpaired electrons calculated from the intensities of the cwEPR spectra, considering the non-central position of the cell in the resonator. c) slices of the subplot a), two component simulations and the dry film (measured in the center of the resonator, without electrolyte).}
\label{fig:S8}
\end{figure}

\begin{figure}
\includegraphics[width=0.88\textwidth]{Figure_S5}
\caption{Spin count during the re-polymerization of a thick (200 deposition cycles), cleaned p-DiTS cell. The cell initially exhibits single broad line but at moderate potentials the sattelite lines are formed. The formation of the three sattelite lines is shown in detail in Figure \ref{fig:S8_SI}. The three EPR lines in the oxidized cell correspond to the released nitroxide fragments that freely tumble in the electrolyte. During the reduction, the satellite peaks disappear and the line broadens. a): CwEPR spectra during the charge-discharge cycle of the cleaned p-DiTS cell that releases the monomers and polymerizes back. b): number of unpaired electrons calculated from the intensities of the cwEPR spectra, considering the non-central position of the cell in the resonator. c) slices of the subplot a), two component simulations and the dry film (measured in the center of the resonator, without electrolyte).}
\label{fig:S8_SI}
\end{figure}

\begin{figure}[h!]
\includegraphics[width=0.5\textwidth]{Figure_S6}
\caption{pDiTS operando redox at RT. Re-polymerisation. Two component fits. Component decomposition. No nitroxide fragments in the used electrolyte.}
\label{fig:S8_SI}
\end{figure}


\begin{figure*}[h!]
\includegraphics[width=0.75\textwidth]{Figure_S7}
\caption{pDiTS operando redox at RT. Fits. Component decomposition. BG.}
\label{fig:S2}
\end{figure*}


\begin{figure*}[h!]
\includegraphics[width=0.5\textwidth]{Figure_S8}
\caption{Initial release of DiTS from an untreated film to the electrolyte.}
\label{fig:Figure3_SI}
\end{figure*}



\begin{figure*}[h!]
\includegraphics[width=0.7\textwidth]{Figure_S9}
\caption{Potential-dependent ex situ cwEPR measurements (X-Band, 9.4$\,$GHz) on a dry, frozen p-DiTS film at 150K. Decomposition of spectral lines into two components with numeric simulations: a broad lorentzian line (g = 2.0070, lw = 4.63$\pm$1.00$\,$mT) and a dilute component with a hyperfine coupling to $^{14}$N (I = 1, g = [2.00231 2.00711 2.00912], A = [100 15 15]$\pm$[1 5 5]$\,$MHz, Lorentzian broadening with lw = 1.65$\,$mT). The absorption profiles are shown with the integrals of the EPR signals. The ratio between the simulated components for each potential is shown by the double integrals of the components. Inset: sum of the broad and dilute components that reproduce the line shape at 0.6$\,$V.\\}
% Spectral lines were deconvoluted into two components by means of numeric simulations with the Easyspin\cite{easyspin} package in Matlab.
%\textbf{DiTS}: dry film /95 cls @ 10 mV/s clean PC+AN / 22dB=1.261915e+00 mW, 5G\\
%\textbf{NiMeOS}: dry film / 15 cls @ 25 mV/s clean AN / 22dB 5G\\
%sim: p-DiTS: pepper, 100\% lorentzian, g=2.00491, lw = 4.63$\,$mT\\
%sim: p-NiMeOS: pepper, $\hat{g}$ and $\Delta \hat{g}$.}
% old: /net/grouphome/ag-bittl/EPR_on_Batteries/Figure_preparation/Figure5/Figure5_params_cleaned_mods
\label{fig:S4}
\end{figure*}




\begin{figure*}[h!]
\includegraphics[width=0.5\textwidth]{Figure_S10}
\caption{NiMeOS ex situ redox at 150K. Fits with $\hat{g}$ and $\hat{g}_{strain}$. }
\label{fig:S3}
\end{figure*}


\begin{figure*}[h!]
\includegraphics[width=0.6\textwidth]{Figure_S11}
\caption{pDiTS operando redox at RT. Release of monomer. Two component fits. Component decomposition. Nitroxide fragments in the used electrolyte.}
\label{fig:S7}
\end{figure*}



\begin{figure*}[h!]
\includegraphics[width=0.3\textwidth]{Figure_S12}
\caption{Top four spectra: release of the DiTS monomer from an untreated pDiTS film to the electrolyte (0.1$\,$M solution of Et$_4$NBF$_4$ in Acetonitrile). Field modulation amplitude: 0.10$\,$mT (overmodulated). Bottom: spectrum of DiTS monomer in the electrolyte at a concentration of 100~$\mu$M for comparison. Field modulation amplitude: 0.02$\,$mT.}
\label{fig:S7_SI}
\end{figure*}


\begin{figure*}[h!]
\includegraphics[width=1\textwidth]{Figure_S13}
\caption{a) Evolution and component decomposition of cwEPR spectra for a tube-based, 150~CV p-DiTS ORB upon 36 charge-discharge cycles. b) Quantitative analysis of the separated spectral components. $\nu$: 9.4$\,$GHz Mod. Amp: 0.5$\,$mT. Quantitative analysis of the CV curves. c) Cyclic voltammograms of the cell, recorded with respect to Ag/AgCl RE at a rate of 5 mV/s inbetween the EPR scans, in the modified tube inside the microwave resonator. The CV were taken at 5$\,$mV/s, which is 10 times slower than the CV curves taken in the larger beaker, because limited amount of the electrolyte in the modified tube was found to strongly affect the shape and positions of the CV peaks at higher speeds of cycling.}
\label{fig:S_Deg}
\end{figure*}

\begin{figure*}[h!]
\includegraphics[width=1\textwidth]{Figure_S14}
\caption{CV curves vs. time for the charge-discharge cycling of a p-DiTS cell. At the scan rate of 5$\,$mV/s each degradation cycle took 200$\,$s. The first CV curve was recorded in a narrower range of potentials. The second CV curve is not shown.}
\label{fig:S_CV_vs_time}
\end{figure*}

\begin{figure*}[h!]
\includegraphics[width=1\textwidth]{Figure_S15}
\caption{Degradation of capacity of a p-DiTS film upon 36 charge-discharge cycles. Transferred charge calculated as integrals of the reduction branches of the CV curves. By the last charge-discharge cycle the film can accept 10 times less the charge as in the beginning of the cycling.}
\label{fig:S_Q_to_Spins}
\end{figure*}

\begin{figure*}[h!]
\includegraphics[width=1\textwidth]{Figure_S16}
\caption{EPR spectra and spin count for the p-DiTS cell (with electrolyte) before and after the degradation series consisting of 36 charge-discharge cycles. EPR spectrum of the degraded cell with the replaced electrolyte does not show any significant contribution of the three-line component, which indicates that the reduced nitroxide fragments have been released to the electrolyte upon the degradation series.}
\label{fig:S_wet_degradation_spin_count}
\end{figure*}

\begin{figure*}[!h]
\includegraphics[width=1\textwidth]{Figure_S17}
\caption{EPR spectra and spin count for the dry p-DiTS film before and after the degradation series consisting of 36 charge-discharge cycles. The reduction of the EPR intensity corresponds to a loss of 1.82e+15 paramagnetic species upon the degradation series.}
\label{fig:S_dry_degradation_spin_count}
\end{figure*}

\subsection{Affection of the Q factor by spectrometer parameters}
\ds{When switching the measurement mode of the X-band spectrometer from 'operate' to 'tune' a shift of the detectors bias current from its ideal value of $200\,\mu$A may occur. Additionally, this shift depends on the microwave power, even when the resonator is critically coupled. To investigate the influence on the Q factor, the tunpicture of the resonator was recorded and analysed for a sample with low absorption - TEMPO in Toluene - and for pDiTS in acetonitrile-based electrolyte, which has a high absorption. While the TEMPO sample is stored in a standard EPR tube and centered in the resonator, the pDiTS sample is put into a modified tube and put at a height of $12\,$mm from the center of the resonator. The tunepicture was recorded for two different dip positions with slight frequency correction, to have a better chance to fit the data. The frequency during the whole measurement never shifted more then $0.5\,\%$.}
\bigbreak
\ds{For the Q factor measurements at different attenuations, two cases were considered: setting the bias current back to $200\,\mu$A after switching to 'tune' mode, which will be referred to as bias adjusted, and not setting the current after the mode switch, and leave the bias not adjusted. Important to note is, that for measuring the pDiTS sample without adjustments, the bias was also not corrected during the change of attenuations. Since this sample cannot be critically coupled, the resulting bias current shift when changing the mode is even bigger than in the TEMPO sample.}
\bigbreak
\begin{figure*}[!ht]
\includegraphics[width=1\textwidth]{Figure_S18}
\caption{\ds{Q factor of the Lyra resonator for TEMPO in Toluene (blue axis) and pDiTS in electrolyte (red axis) from fitting the tunepicture for different microwave attenuations, with and without adjusting the diode bias current to $200\,\mu$A}}
\label{fig:S_LyraQ_power}
\end{figure*}
\ds{
The results in Figure \ref{fig:S_LyraQ_power} show, that with adjustment of the bias current, both samples show a similar trend for the Q factor, namely a constant region in the center of the recording range and decreasing Q factors when moving away from the center. For pDiTS, this trend shows a bigger relative difference in both directions compared to TEMPO. Without adjusting the bias current, the Q factor stays rougly constant for TEMPO for lower attenuations and follows its behavior as described before for higher attenuations. In contrast, pDiTS starts at a higher Q factor, when the bias is not adjusted and decreases monotonously with increasing attenuation, but this trend flattens towards the end. Within the region of interest for pDiTS, namely $23\,$dB to $35\,$dB, the Q factor is constant within a range of $20\,\%$ of the mean value, while the Q factor at $23\,$dB is $15.5\,\%$ higher than the Q factor at $35\,$dB. This lies within the error considerations of the spincounting procedure. If the bias is not adjusted, the Q factor at $23\,$dB is approximately 2.2 times higher than at $35\,$dB. When comparing the data for adjusted and umadjusted bias current for $23\,$dB, the difference of the Q factor is $18.3\,\%$, relative to the lower value, which is again within the error considerations. At $35\,$dB, this difference is $61.5\,\%$. While it is reasonable to compare quantitative data aquired at $23\,$dB. regardless of the adjustment, it is only possible to compare to data aquired at $35\,$dB, if the bias was adjusted.}
\bigbreak
\ds{A similar measurement series was performed for both samples at constant microwave powers but with different fixed bias currents for aquiring the tunepicture. The results for that are shown in Figure \ref{fig:S_LyraQ_bias}. For both samples again similar trends are observed with the pDiTS data scattering more in general. While for intermediate and high bias currents the Q factors are rather constant, the Q factor drops rapidly when going to small currents. For even lower current values, the tunepicture can not be analysed anymore, because the dip is clipped by the detector. When changing the bias current by $\pm 50\,\mu$A from $200\,\mu$A, the Q factor changes for pDiTS at most $25\,\%$ for all the observed attenuations. But this maximum deviation is only reached for $26\,$dB and $35\,$dB. For $23\,$dB the deviation lies in between $10\,\%$ and for $30\,$dB, it is at most $13\,\%$. This suggests, that the changes of the bias current have to be properly monitored especially for higher attenuations. Because the coupling for pDiTS was optimised for $23\,$dB, the initial shift of bias current at that power is small and the Q factor can be considered stable within the error margin of the spincounting. The results from this series also partly explain the results of the unadjusted attenuation dependence for pDiTS from before, since the initial bias current drops with higher attenuation and small currents lead to a relevant decrease in Q factor.}
\begin{figure*}[!ht]
\includegraphics[width=1\textwidth]{Figure_S19}
\caption{\ds{Q factor of the Lyra resonator for TEMPO in Toluene (blue axis) and pDiTS in electrolyte (red axis) from fitting the tunepicture for different set diode bias current values}}
\label{fig:S_LyraQ_bias}
\end{figure*}

\subsection{Materials}
\subsubsection{Polymers}
Di-TEMPO-Salen and NiMeOSalen were synthesized at the department of electrochemistry, St. Petersburg State University.
\subsubsection{Solvents}
Anhydrous Acetonitrile (ACN) was obtained from the Sigma Aldrich company.
Gold (Au) and Chromium (Cr) pillets for the thermal evaporation were obtained from the Kurt J. Lesker company.
Silver (Ag) wire was obtained from ?


\end{document}
